\section{Casi d'uso implementati}
In questa fase sono stati sviluppati i seguenti casi d'uso ritenuti prioritari per lo sviluppo di OptiTour.

\begin{table}[hbt!]
    \centering
    \begin{tabular}{ll}
        \hline
        \textbf{ID} & \textbf{Titolo} \\
        \hline
        UC1.1 & Registrazione \\
        UC1.2 & Login \\
        UC1.3 & Logout \\
        UC11 & Imposta punto di partenza \\
        UC12.1 & Creazione itinerario manuale \\
        UC13 & Scelta tempistiche \\
        UC3 & Visualizza percorsi \\
        UC16 & Elimina viaggio \\
        UC14 & Ottimizzazione percorso \\
        \hline
    \end{tabular}
    \caption{Iterazione 1}
\end{table}

\subsection*{UC1.1: Registrazione}
\textbf{Attori:} Utente, Sistema. \\
\textbf{Descrizione:} Un nuovo utente può creare un account personale fornendo i dati richiesti. \\
\textbf{Flusso degli eventi:}
\begin{enumerate}
    \setlength{\itemsep}{0pt}
    \setlength{\parskip}{0pt}
    \item L'utente accede alla pagina di registrazione.
    \item Inserisce i propri dati personali e le credenziali.
    \item Il sistema verifica i dati e crea il nuovo account.
    \item L'utente viene reindirizzato alla pagina di login.
\end{enumerate}

\subsection*{UC1.2: Login}
\textbf{Attori:} Utente, Sistema. \\
\textbf{Descrizione:} Un utente registrato può accedere al proprio account inserendo le proprie credenziali. \\
\textbf{Flusso degli eventi:}
\begin{enumerate}
    \setlength{\itemsep}{0pt}
    \setlength{\parskip}{0pt}
    \item L'utente accede alla pagina di login.
    \item Inserisce email e password.
    \item Il sistema verifica le credenziali e autentica l'utente.
    \item L'utente viene reindirizzato alla home.
\end{enumerate}

\subsection*{UC1.3: Logout}
\textbf{Attori:} Utente, Sistema. \\
\textbf{Descrizione:} L'utente può terminare la propria sessione in modo sicuro. \\
\textbf{Flusso degli eventi:}
\begin{enumerate}
    \setlength{\itemsep}{0pt}
    \setlength{\parskip}{0pt}
    \item L'utente clicca sul pulsante di logout.
    \item Il sistema termina la sessione corrente.
    \item L'utente viene reindirizzato alla pagina di login.
\end{enumerate}

\subsection*{UC11: Imposta punto di partenza}
\textbf{Attori:} Utente, Sistema \\
\textbf{Descrizione:} L'utente inserisce un indirizzo testuale come punto di partenza del viaggio. Il sistema converte automaticamente l'indirizzo in coordinate geografiche tramite il servizio esterno Nominatim e le salva nel viaggio. \\
\textbf{Flusso degli eventi:}
\begin{enumerate}
    \setlength{\itemsep}{0pt}
    \setlength{\parskip}{0pt}
    \item L'utente inserisce un indirizzo di partenza nel campo apposito.
    \item Il sistema invia l'indirizzo al servizio Nominatim.
    \item Nominatim restituisce le coordinate geografiche corrispondenti.
    \item Le coordinate vengono salvate nel viaggio e la mappa si aggiorna sul punto selezionato.
\end{enumerate}

\subsection*{UC12.1: Creazione itinerario manuale}
\textbf{Attori:} Utente, Sistema. \\
\textbf{Descrizione:} L'utente seleziona manualmente i monumenti che desidera visitare. I monumenti vengono recuperati da Overpass API la prima volta e salvati in MongoDB per evitare chiamate ripetute al servizio esterno. \\
\textbf{Flusso degli eventi:}
\begin{enumerate}
    \setlength{\itemsep}{0pt}
    \setlength{\parskip}{0pt}
    \item L'utente cerca i monumenti di una città.
    \item Il sistema verifica se i monumenti sono già presenti in MongoDB.
    \item Se non presenti, il sistema li recupera tramite Overpass API e li salva nel database.
    \item I monumenti vengono mostrati sulla mappa.
    \item L'utente seleziona i monumenti che desidera visitare.
    \item L'utente conferma la selezione e il sistema salva il viaggio.
\end{enumerate}

\subsection*{UC13: Scelta tempistiche}
\textbf{Attori:} Utente, Sistema. \\
\textbf{Descrizione:} L'utente può definire per ciascun monumento selezionato il tempo che intende dedicare alla visita, espresso in minuti. \\
\textbf{Flusso degli eventi:}
\begin{enumerate}
    \setlength{\itemsep}{0pt}
    \setlength{\parskip}{0pt}
    \item L'utente seleziona un monumento durante la creazione del viaggio.
    \item Inserisce il tempo di visita desiderato in minuti.
    \item Il sistema salva il tempo di visita associato a quella tappa del viaggio.
\end{enumerate}

\subsection*{UC3.2: Visualizza percorsi}
\textbf{Attori:} Utente, Sistema. \\
\textbf{Descrizione:} L'utente può visualizzare i propri viaggi e di visualizzare i dettagli del singolo viaggio. \\
\textbf{Flusso degli eventi:}
\begin{enumerate}
    \setlength{\itemsep}{0pt}
    \setlength{\parskip}{0pt}
    \item L'utente accede alla sezione dei propri viaggi.
    \item Il sistema recupera tutti i viaggi dell'utente dal database.
    \item L'utente può selezionare un singolo viaggio per visualizzarne i dettagli.
\end{enumerate}

\subsection*{UC16: Elimina viaggio}
\textbf{Attori:} Utente, Sistema. \\
\textbf{Descrizione:} L'utente può eliminare un viaggio esistente. \\
\textbf{Flusso degli eventi:}
\begin{enumerate}
    \setlength{\itemsep}{0pt}
    \setlength{\parskip}{0pt}
    \item L'utente seleziona il viaggio che desidera eliminare.
    \item Conferma l'operazione di eliminazione.
    \item Il sistema rimuove il viaggio dal database.
\end{enumerate}

\subsection*{UC14: Ottimizzazione percorso}
\textbf{Attori:} Utente, Sistema \\
\textbf{Descrizione:} Il sistema ottimizza l'ordine delle tappe del viaggio, calcolando il percorso più efficiente tra i monumenti selezionati a partire dal punto di partenza definito dall'utente. \\
\textbf{Flusso degli eventi:}
\begin{enumerate}
    \setlength{\itemsep}{0pt}
    \setlength{\parskip}{0pt}
    \item Il sistema recupera il viaggio tramite l'ID.
    \item Legge le coordinate del punto di partenza e di ogni monumento dal database.
    \item Applica l'algoritmo di ottimizzazione per calcolare il percorso ottimale.
    \item Il viaggio viene aggiornato con il nuovo ordine delle tappe e lo stato viene impostato a SAVED.
\end{enumerate}